\section{Weighted cactus optimization in accuracy bits}
\label{sec:a-weighted-cactus}
\label{sec:a-open-problems}

A cactus may have arbitrarily many cycle blocks. Even when its nominations
depend on one scalar, a weighted potential objective can therefore contain
arbitrarily many independent radicals. The algorithms in this section
approximate the local formulas before adding them, so that the final
semialgebraic optimization has only the nomination parameters as variables.
A structural theorem reduces an unrestricted balanced nomination box
to polynomially many such parameter families when the objective support
and maximum block rank are fixed.

Throughout, $c\in\Q^V$ satisfies $\one^\top c=0$ and the objective is
$W=c^\top\pi$. Accuracy-bit complexity means polynomial dependence on the
binary input length and $\max\{0,\log(1/\eps)\}$ for positive rational
$\eps$; we may replace $\eps$ by $\min\{\eps,1\}$. The exponents of our
polynomial bounds may depend on the fixed parameters, so no
fixed-parameter tractability is claimed.

\begin{theorem}[Affine nominations on a cactus]
\label{thm:a-wcac-affine}
Fix $d$. Let $G$ be a connected cactus and let
$b(z)=b^0+Hz$ be a rational affine balanced nomination family on a nonempty
compact rational polytope $P\subseteq\R^d$. The support of $c$, cycle
lengths, and number of cycles are unrestricted. Consider either
\begin{enumerate}
\item fixed positive rational asymmetric quadratic edge laws; or
\item uncertain quadratic laws in either of two models: one shared
coefficient $\beta_e\in[L_e,U_e]$ per symmetric edge, or two independent
directional coefficients $\beta_e^\pm\in[L_e^\pm,U_e^\pm]$ per asymmetric
edge. All intervals have positive rational endpoints and are independent
across edges.
\end{enumerate}
For each maximum and minimum of $W$, a rational interval of width at most
$\eps$ containing the optimum and a rational $\eps$-optimal original
parameter scenario can be computed in accuracy-bit polynomial time.
In case~2 the extrema are joint over $z$ and all resistances. No operating
constraints filter the scenario set in this statement.
\end{theorem}

The local quadratic and radical formulas used below have direct
predecessors. Gotzes, Heitsch, Henrion, and Schultz derive piecewise
quadratic-radical and rational circulation formulas along affine
nomination rays in Theorem~6 of the author manuscript of
\citet{gotzes2016-quantification}; that manuscript also extends the
method to node-disjoint cycles with attached trees (printed pp.~18
and~24). Vigneron optimizes sums of algebraic functions by common
arrangements and approximate summation
\citep[Theorem~6 and Section~3.2 of the October~21, 2011 author manuscript]{vigneron2014-geometric-optimization-and-sums-of}.
His Section~2.3 includes a bit model, in which the approximation schemes
are polynomial in $1/\eps$ rather than in accuracy bits. Thus neither the
local formulas nor the use of approximation to avoid exact radical-sum
comparisons is a new principle here, and we make no claim of exhaustive
literature priority. The guarantee established below combines arbitrary
weighted objectives, fixed nomination dimension, unbounded cactus cycle
count, the stated law models, explicit accuracy-bit bounds, and parameter
recovery.

\subsection{A resistance-independent nomination-face theorem}

A \emph{nomination face} fixes some coordinates at their box endpoints,
lets the others keep their intervals, and imposes exact total balance.
Faces in the next theorem are formed after the stated
objective-preserving reductions, so they need not be faces of the original
box.

\begin{theorem}[Small nomination faces at bounded block rank]
\label{thm:a-wcac-faces}
Fix integers $r\ge1$ and $p\ge2$. Suppose every block of a connected graph
has cycle rank at most $r$, the nominations range over a nonempty rational
balanced box, and $|\supp c|\le p$. For fixed positive quadratic laws,
including asymmetric laws, one can construct $n^{O(rp)}$ rational faces,
each with $O(rp)$ unfixed nomination coordinates, such that a maximum of
$W$ lies on at least one face after objective-preserving reductions.
Every feasible reduced nomination has a rational lift when it is rational,
with polynomial encoding length.

The reductions and face family depend on the graph, nomination bounds,
and $c$, but not on the positive law coefficients. Consequently the same
family contains an optimum for joint optimization over any independent
compact positive coefficient sets. This last assertion is structural;
it does not compute the variable-coefficient block objectives.
\end{theorem}
\begin{proof}
If $c=0$, start at all nomination lower bounds and raise coordinates in a
fixed order until their sum is zero; feasibility ensures that this is
possible, and all but at most one coordinate end at endpoints. Fixing
those endpoint coordinates and retaining the possible remaining interval
with exact balance gives one rational face with at most one unfixed
coordinate, using rational arithmetic of polynomial encoding length.

Otherwise use the block--vertex incidence tree, with a node for every
block and every original vertex. A subtree with no objective support
except possibly at its single attachment can be pruned by adding all its
nomination intervals to the attachment interval: its exterior state
depends only on that aggregate, and its removed vertices have zero
objective coefficient. A prescribed aggregate in a sum of intervals can
be split by starting at all lower endpoints and filling the remaining
amount in a fixed order, with rational arithmetic of polynomial encoding
length.

A second reduction is needed. For a block $B$, remove its edges and
write $G_v$ for the component rooted at $v\in V(B)$. Put
$\gamma_v=\sum_{u\in V(G_v)}c_u$. If every $\gamma_v=0$, contract all of
$B$ to one vertex, summing its vertices' nomination intervals and objective
coefficients. To prove that the objective is preserved, write
\begin{equation}
 W=\sum_{v\in V(B)}\sum_{u\in V(G_v)}c_u(\pi_u-\pi_v)
       +\sum_{v\in V(B)}\gamma_v\pi_v.
 \label{eq:a-wcac-contraction}
\end{equation}
The last sum vanishes. The flows and relative potentials in each $G_v$
are determined by the nominations away from its root, and the root
supplies their balancing nomination. Translating these components to a
common root potential and summing their root nominations produces the
contracted state and leaves \eqref{eq:a-wcac-contraction} unchanged.
Conversely, disaggregate the contracted nomination within the original
root intervals. Subtracting the already determined exterior flows leaves
balanced effective injections into $B$. Solve the passive block and
translate each $G_v$ to its recovered root potential. Existence and
uniqueness follow from \cref{thm:a-pre-existence}, and the objective is
again unchanged. Any admissible coefficients on the contracted block
suffice for this lift.

Repeat pruning and zero-$\gamma$ contraction until neither applies.
No surviving block loses two distinct vertices to one contraction, since
that would create a cycle in the incidence tree. Thus surviving block
ranks are unchanged, and objective support cannot increase. The induced
coefficient sums on surviving blocks are unaffected, because contraction
only regroups terms in the same components. All tests use rational sums
of $c$ and do not inspect resistances. The reductions use only nomination
boxes and balance; flow or potential filters would invalidate them.

Let $T$ be the minimal incidence subtree spanning the remaining support.
Retain each physical block represented in $T$ in its entirety, including
its nonsupport vertices and their nominations. Every leaf of $T$ is a
support vertex. A tree with at most $p$ leaves has at most $p-2$ branching
nodes and total branching degree $O(p)$. Mark the support vertex nodes
and branching vertex nodes. Call a block \emph{special} if its block node
branches in $T$ or is adjacent to a marked vertex node. There are $O(p)$
special blocks, with $O(p)$ total incidences in $T$, and all other blocks
form $O(p)$ \emph{ordinary chains}. Each ordinary block has two retained
ports and no internal objective source or retained side attachment. Its
induced sources for an electrical adjoint are $I,-I$, with the same $I$
along the chain in a consistent orientation. This $I$ is a known sum of
objective coefficients, and it is nonzero because otherwise the block
would have been contracted.

First smooth the laws to $g_e^\rho(t)=g_e(t)+\rho t$, $\rho>0$, whose
differential resistances are positive, also at zero flow.
Implicit differentiation as in \cref{lem:a-cac-adjoint} gives the
nomination gradient $h$, modulo a constant, from
$A\operatorname{diag}(1/(g_e^\rho)'(x_e))A^\top h=c$.
By the two-terminal maximum principle, the internal $h$ values of an
ordinary nonbridge block lie strictly between its port values. Indeed an
internal extremum propagates through all neighbors unless it reaches a
terminal, and two-vertex-connectivity gives a path to the other terminal
avoiding the first. The two terminal values differ because the injected
current $I$ is nonzero. Bridge blocks have no internal vertices.
Along each chain the port values are strictly ordered, in the direction
fixed by $I$, and consecutive open internal ranges are disjoint.

At a maximum on the balanced box the polyhedral normal-cone condition
supplies a scalar $\lambda$ such that
\begin{equation}
 h_v>\lambda\ \Longrightarrow\ b_v=u_v,
 \qquad h_v<\lambda\ \Longrightarrow\ b_v=\ell_v.
 \label{eq:a-wcac-kkt}
\end{equation}
This condition remains valid for singleton intervals and a
lower-dimensional balanced box. On each ordinary chain, at most one
block's open range contains $\lambda$. If $\lambda$ equals an articulation
level, choose either adjacent block and leave that block's endpoints free.
If $\lambda$ is outside the chain range, all internal coordinates saturate
on one side. Enumerate a selected block on each chain, or one of the two
all-saturated alternatives. Outside the selected blocks the endpoint
choices are fixed by the known sign of $I$; all vertices of special blocks
remain free. A shared endpoint is left free whenever a selected or special
block contains it, so no inconsistent assignment arises. This gives
$n^{O(p)}$ closed \emph{coarse faces}, each involving only $O(p)$ free
blocks.

We next justify taking limits before the second perturbation.
All physical flows are bounded by a common total-positive-nomination
bound $B_0$; for $0<\rho\le1$, edge drops are bounded by
$\beta_U B_0^2+B_0$, and path sums bound normalized potentials. States are
jointly continuous in nominations and $\rho\ge0$: a fixed spanning-tree
right inverse of $A$ corrects a feasible flow when nominations vary, and
energy minimality and strict convexity identify every subsequential limit
as the unique limiting physical flow. Edge laws and path integration then
give convergence of potentials, and compactness gives uniform convergence
on the nomination box as $\rho\downarrow0$. A subsequence of smoothed
maximizers belongs to one fixed closed coarse face, and its limit is an
original maximizer. Fix this coarse face.

Within its special and selected nonbridge blocks, mark every retained
port, every objective-support vertex, and every vertex of internal block
degree at least three. A block of rank $r_B$ has at most $2r_B-2$ vertices
of the last kind, since $\sum_v(\deg_Bv-2)=2r_B-2$. The total port count
is $O(p)$ and there are $O(p)$ such blocks, so $O(rp)$ vertices are
marked. Suppress unmarked degree-two interiors into paths. By Euler's
formula, a block with $k_B$ marked vertices has $r_B+k_B-1$ such paths;
this includes a cycle with two marked ports. Every retained block has at
least two ports. Mark bridge endpoints directly. In total there are
$O(rp)$ paths and marked coordinates.

Smooth again and, only on the fixed coarse face, perturb the objective by
\[
 \delta\sum_{v\in U}(\pi_v-\pi_a),\qquad \delta>0,
\]
where $U$ consists of these path-interior vertices and $a$ is marked.
Vertices of $U$ have no original objective sources or off-block
attachments, so each has adjoint source exactly $\delta$.
Orient a suppressed path from left to right and let $j_i$ be its
successive adjoint currents. Conservation gives
$j_i-j_{i-1}=\delta$. Because voltage decreases in the direction of positive
current, the adjoint values first increase and then decrease, with at
most a plateau at two consecutive vertices, so a horizontal level contains
at most two interior vertices. Applying \eqref{eq:a-wcac-kkt} on this
coarse face shows that each path has lower endpoints, at most one free
pivot, upper endpoints, at most one free pivot, and lower endpoints;
empty runs are allowed. Enumerate the two pivot positions and endpoint
runs, and leave marked coordinates free. This gives $n^{O(rp)}$ faces,
each with $O(rp)$ free coordinates.

The perturbation is uniformly at most $\delta|U|$ times the uniform
potential-difference bound. Letting $\delta,\rho\downarrow0$ and passing
to a fixed-face subsequence again preserves a maximum on the coarse face,
hence a global maximum. The coarse face must be fixed before the positive
interior sources are added: perturbing all vertices at the first stage
would destroy the zero-source chain argument.

Every construction and endpoint pattern is independent of the
coefficients. After balance elimination, each face is a compact rational
polytope with $O(rp)$ coordinates and a rational affine nomination map.
Interval sums and their greedy disaggregation have polynomial bit length.
For joint coefficient optimization, choose an attained joint maximum,
freeze its coefficients, and apply the same family. The reductions preserve
attainable objective values at those surviving coefficients, and arbitrary
admissible removed coefficients permit a lift. This proves the final
assertion.
\end{proof}

\subsection{Fixed-law cycle formulas and rational square-root panels}

Choose a rational spanning-tree flow $w$ with $Aw=c$, putting zero on
chords. Then
\begin{equation}
 W=w^\top A^\top\pi=\sum_e w_eg_e(x_e).
 \label{eq:a-wcac-weights}
\end{equation}
A spanning-tree particular flow for $b(z)$ is affine in $z$. Coherently
orient each cycle, reversing its offsets and weights and swapping
asymmetric coefficients and their directional intervals when necessary.
On a cycle its flows have form
\begin{equation}
 x_e=q+\ell_e(z),\qquad
 F(z,q)=\sum_{e\in C}g_e(q+\ell_e(z))=0,
 \label{eq:a-wcac-cycle}
\end{equation}
with rational affine $\ell_e$. Bridges have affine flows. The cycle
variables are independent because the cycles are edge-disjoint, and
cycle-compatible drops integrate to a global potential. Each $F(z,\cdot)$
is continuous and strictly increasing, with opposite limits at infinity.

Rational linear programming supplies a box containing $P$ whose endpoint
bit lengths are polynomial, even if $P$ has lower dimension.
Coefficient bounds on this box supply a polynomial-bit $B_0\ge1$ bounding
total positive nominations, hence all physical flows, and
$R\ge1$ with $|q|\le R$ in every physical scenario; if a chord offset is
zero, $R=B_0$ suffices. Write $\beta_U$ for the largest coefficient or
allowed coefficient endpoint. The uniform bound
\begin{equation}
 |W|\le \|c\|_1(n-1)\beta_U B_0^2
 \label{eq:a-wcac-value-bound}
\end{equation}
follows by normalizing one potential and integrating drops along paths.

\begin{lemma}[Local fixed-law representation]
\label{lem:a-wcac-local}
For fixed $d$, a cycle's contribution to \eqref{eq:a-wcac-weights} has
polynomially many exact semialgebraic charts in $z$, each represented by
a rational function of bounded degree or by
\[
 P_2(z)+L_1(z)\sqrt{\Delta(z)},
\]
where $P_2,\Delta$ are rational polynomials of degree at most two and
$L_1$ is rational affine. Chart domains have polynomial-size
bounded-degree rational descriptions. All bit lengths are polynomial.
The all-zero-flow stratum is included explicitly.
\end{lemma}
\begin{proof}
Enumerate the realizable sign cells of the affine forms $q+\ell_e(z)$ in
$d+1$ variables, and take their closures. For fixed $d$ there are
polynomially many, so all $2^{|C|}$ sign vectors need not be enumerated.
On each closed cell, \eqref{eq:a-wcac-cycle} is
\[
 A_2q^2+B_1(z)q+C_2(z)=0,
\]
where $A_2$ is rational constant and the subscripts bound degrees. At a
physical solution its derivative satisfies
\[
 2A_2q+B_1(z)=2\sum_{e\in C}\beta_e^{\sgn x_e}|x_e|\ge0.
\]
For $A_2\ne0$ the physical root is therefore
\begin{equation}
 q=\frac{-B_1(z)+\sqrt{B_1(z)^2-4A_2C_2(z)}}{2A_2}.
 \label{eq:a-wcac-root}
\end{equation}
The plus sign is correct even when $A_2<0$. Substituting into a quadratic
weighted drop gives the asserted polynomial plus affine radical term;
all divisions are by the fixed nonzero rational $2A_2$ or its square.
For $A_2=0$ and $B_1(z)>0$, substitute $q=-C_2/B_1$ exactly and keep the
nonzero-denominator condition; no uniform lower bound on $B_1$ is needed.
If $A_2=B_1(z)=0$ at a physical point, the derivative identity forces every
flow to vanish. Represent this chart by
$q=-\ell_1(z)$ and $\ell_e(z)=\ell_1(z)$ for every $e$; its contribution
is zero. These charts cover all physical solutions.

For each chart, eliminate its single $q$ from the root equation, the
closed sign inequalities, $|q|\le R$, and, when needed, the root selector
$2A_2q+B_1\ge0$. In the linear case keep $B_1>0$; in the zero
case use the displayed affine equalities. The fixed-dimensional
elimination bounds in \eqref{eq:a-pre-qe-cost} give polynomially many
rational polynomials of bounded degree and polynomial coefficient bit
length. Disconnected domains are allowed. By uniqueness, several valid
charts at a boundary give the same physical contribution.
\end{proof}

\begin{lemma}[Constructive rational square-root panels]
\label{lem:a-wcac-panels}
Given a rational $S\ge1$ of polynomial bit length and $0<\eta\le1$, one
can construct a piecewise rational polynomial approximation to $\sqrt t$
on $[0,S^2]$, uniformly within $\eta$, whose panel count, degrees, dense
coefficient sizes, and construction time are polynomial in the input
length and $\log(1/\eta)$.
\end{lemma}
\begin{proof}
First take $S=1$. Choose $J$ with $2^{-J}\le\eta$ and use zero on
$[0,4^{-J}]$. For $j=0,\ldots,J-1$, use the panel
$[4^{-j-1},4^{-j}]$, center $a_j=(9/16)4^{-j}$, and polynomial
\[
 T_j(t)=\frac34\,2^{-j}
       \sum_{k=0}^K \binom{1/2}{k}(t/a_j-1)^k.
\]
All numbers are rational. On the panel $|t/a_j-1|\le7/9$ and the binomial
coefficients have magnitude at most one. The convergent binomial series
therefore gives
\[
 |T_j(t)-\sqrt t|\le\frac92(7/9)^{K+1}.
\]
Choose $K=O(1+\log(1/\eta))$ to make this at most $\eta$. On the zero
panel the error is at most $2^{-J}\le\eta$, the value of $\sqrt t$ at its
right endpoint. Either adjacent formula is valid at a shared boundary;
the approximant need not be continuous. Expanded coefficient bit lengths
are polynomial in $J,K$, including the products $jK$. For general $S$,
approximate $\sqrt{t/S^2}$ within $\eta/S$ and multiply by $S$. This
proves all bounds.
\end{proof}

For each discriminant in \cref{lem:a-wcac-local}, absolute coefficient
sums on the box for $P$ give a rational power of two $S$ of polynomial
bit length bounding its square root on its domain. Absolute coefficient
bounds similarly give $M_i\ge |L_{1,i}(z)|$. Set
\begin{equation}
 \eta=\frac{\eps}{16(1+\sum_i M_i)},\qquad a=\eps/16,
 \label{eq:a-wcac-error}
\end{equation}
where the sum may range over all charts. Replacing every radical by its
panels changes any selected sum of local contributions by at most $a$,
and all rational charts are kept exact. Substituting a quadratic
discriminant into a panel polynomial doubles its degree; a dense
polynomial of polynomial degree in fixed $d$ variables still has
polynomially many monomials.

\subsection{One common partition and original-parameter recovery}

Collect the polynomials appearing in every chart-domain formula, every
panel-boundary formula $\Delta-S^2t_j$, every rational denominator, the
bridge zero-flow hyperplanes, and the inequalities for $P$. Fixed-dimensional
sign decomposition, or sampling all realizable sign conditions followed
by their Boolean descriptions, gives a polynomial-size common partition.
On every part choose an eligible chart and panel for each cycle, and the
correct bridge-law branch; eligibility depends only on these signs.
Thus every point is covered, including all lower-dimensional boundaries
and singular zero-flow strata, and combinations of charts across cycles
are never enumerated.

The sum of the chosen surrogates is one rational function $\widehat W(z)$
on each part. Multiplying polynomially many bounded-degree denominators
and adding numerators gives polynomial degree and polynomial coefficient
bit length, and dense expansion takes polynomial time since $d$ is fixed.
For threshold decisions one may multiply by the squared common
denominator, retaining all nonzero conditions. The real-algebraic bounds
from the preliminaries therefore apply, even though the degrees now grow
with accuracy bits.

The parts may be open. We therefore optimize a supremum with the
existential test $\exists z\text{ in the part}:\widehat W(z)>t$, without
assuming attainment on every part. The uniform error $a$ and
\eqref{eq:a-wcac-value-bound} give a polynomial-bit bound on every
surrogate supremum, including near a vanishing chart denominator.
Bisection on the union of these tests produces an interval $[L,U]$ of
width at most $\eps/8$ containing the largest surrogate supremum.
Then
\begin{equation}
 [L-a,U+a]
 \label{eq:a-wcac-certified-interval}
\end{equation}
contains the physical optimum and has width at most $\eps/4$.
Sample on a part at the strictly lower threshold $L-\eps/16$;
the existential sampling primitive \cref{prim:a-pre-sample-points}
returns an algebraic $z_*$ of polynomial representation size. Its true
objective is within $5\eps/16$ of the optimum, since the optimum is at
most $U+a$ and the sample has value greater than $L-\eps/16-a$. The
strictly lower threshold handles unattained suprema and exact equality
alike.

\begin{lemma}[Nomination recovery bound]
\label{lem:a-wcac-nomination-continuity}
For fixed asymmetric quadratic laws and two balanced nominations whose
physical flows have magnitude at most $B_0$,
\begin{align}
 \|x(b)-x(b')\|_\infty&\le\tfrac12\|b-b'\|_1,
 \label{eq:a-wcac-flow-nomination}\\
 |W(b)-W(b')|&\le K_b\|b-b'\|_1,
 \qquad K_b=\|c\|_1(n-1)\beta_U B_0.
 \label{eq:a-wcac-nomination-bound}
\end{align}
These inequalities hold on every connected graph.
\end{lemma}
\begin{proof}
Put $y=x(b)-x(b')$ and $h=\pi(b)-\pi(b')$. Strict monotonicity gives
$\sgn y_e=\sgn(A^\top h)_e$ whenever $y_e\ne0$. Orient these difference
flows positively. The potential $h$ strictly decreases along each such
arc, so their directed support is acyclic. A path decomposition from
positive to negative entries of $b-b'$ has total weight $\|b-b'\|_1/2$,
which proves the first inequality. On $[-B_0,B_0]$ each law is Lipschitz
with constant $2\beta_U B_0$, also across zero, so every drop changes by
at most $\beta_U B_0\|b-b'\|_1$. Normalize one potential, use paths of
length at most $n-1$, and sum against $|c_v|$.
\end{proof}

Let $H_1=\sum_{v,j}|H_{vj}|$. Isolate each coordinate of $z_*$ in a rational
interval of width $h$ with $K_bH_1h\le\eps/4$; this condition is omitted
if its coefficient is zero. Rational linear programming on the
intersection of that box with $P$ returns a rational $z'\in P$ of
polynomial bit length. The intersection contains $z_*$, even for
lower-dimensional $P$, and $\|b(z')-b(z_*)\|_1\le H_1h$, so the physical
objective loss is at most $\eps/4$. Rounding need not preserve the
selected chart, panel, or any nonzero denominator; only the global
physical estimate matters. At rational $z'$ the cycle roots can be
enclosed individually by monotone bisection and their drops summed with
rational error bounds, without a common number field. This proves case~1
of \cref{thm:a-wcac-affine}; negating $c$ proves the minimum statement.

\subsection{Eliminating independent resistance intervals}

We now prove case~2 of \cref{thm:a-wcac-affine}. Independent positive
resistance intervals are an established gas-network uncertainty model.
A\ss mann, Liers, Stingl, and Vera give exact tree reductions by linear
programming and a single-cycle reduction of circulation intervals to
polyhedral coefficient conditions
\citep[Sections~4.2 and~4.3.3, Lemma~4.10]{amann2018-deciding-robust-feasibility-and-infeasibility}.
Gonz\'alez Grand\'on, Heitsch, and Henrion combine random loads and robust
roughness uncertainty on trees, with explicit rectangular and ellipsoidal
inner optimizations \citep[Sections~2--3]{gonzalez2017-joint}.
The additional task here is deterministic weighted optimization over
independent cycle blocks with the accuracy and recovery guarantees of
\cref{thm:a-wcac-affine}.

For fixed $z$, \eqref{eq:a-wcac-weights} separates into independent cycle
maxima and bridge maxima. Each bridge chooses the endpoint maximizing its
weighted drop, using its active directional interval in the asymmetric
model. On a closed flow-sign cell, write $[L_e,U_e]$ for the active
interval $[L_e^{s_e},U_e^{s_e}]$ in that model; in the symmetric model
retain the shared interval. Only the active coefficient enters the local
problem: inactive coefficients can take any allowed rational value, and
at zero flow neither coefficient contributes. Thus the two models have
the same local reduction, although shared symmetric uncertainty is a
diagonal constraint and not literally a product of two directional
intervals. On a closed flow-sign cell of a cycle, put
$f_e(z,q)=s_e(q+\ell_e(z))^2$, where $s_e\in\{-1,1\}$ is the selected
sign; at zero either sign gives the correct value. For fixed $(z,q)$ the
resistance problem is the one-equality box linear program
\begin{equation}
 \max_{L\le\beta\le U}
 \left\{\sum_e w_ef_e\beta_e:\ \sum_e f_e\beta_e=0\right\}.
 \label{eq:a-wcac-local-lp}
\end{equation}

\begin{lemma}[Threshold elimination with ties]
\label{lem:a-wcac-threshold}
For each distinct edge weight $\lambda\in\{w_e:e\in C\}$ let
$T=\{e:w_e=\lambda\}$. For $e\notin T$, choose $\widehat\beta_e=U_e$
when $(w_e-\lambda)s_e>0$, and $L_e$ otherwise. Define
\begin{align*}
 S&=\sum_{e\notin T}\widehat\beta_ef_e,\\
 l_T&=\sum_{e\in T}\min\{L_ef_e,U_ef_e\},\qquad
 u_T=\sum_{e\in T}\max\{L_ef_e,U_ef_e\}.
\end{align*}
This threshold has a feasible optimizing profile exactly on the domain
\begin{equation}
 S+l_T\le0\le S+u_T.
 \label{eq:a-wcac-tie-domain}
\end{equation}
Its value there is the rational quadratic polynomial
\begin{equation}
 Q_\lambda(z,q)=\sum_{e\notin T}(w_e-\lambda)\widehat\beta_ef_e(z,q).
 \label{eq:a-wcac-threshold-value}
\end{equation}
If \eqref{eq:a-wcac-local-lp} is feasible, at least one threshold applies.
\end{lemma}
\begin{proof}
The tied coordinates fill every aggregate value in $[l_T,u_T]$, so they
can supply $-S$ exactly when \eqref{eq:a-wcac-tie-domain} holds. Every
untied coordinate maximizes its reduced-cost expression
$(w_e-\lambda)f_e\beta_e$, also in the zero-flow case, where either
endpoint is optimal, and the tied reduced costs vanish. The feasible
profile therefore attains the Lagrangian upper bound, whose value is
\eqref{eq:a-wcac-threshold-value} because $\sum f_e\beta_e=0$.

Conversely, the dual function is the convex piecewise-affine function
\[
 D(\lambda)=\sum_e\max_{L_e\le\beta_e\le U_e}
                     (w_e-\lambda)f_e\beta_e.
\]
Partition the real line at all distinct listed edge weights, whether or
not they are actual breakpoints of $D$. On every resulting interval or
tail, $D$ is affine. Feasibility bounds $D$ below by the primal optimum,
so its tails cannot decrease without bound. A minimum therefore occurs
at a listed weight: an affine function on a bounded interval has a
minimizing endpoint, and a flat tail has a minimizing finite endpoint.
This also covers a globally constant dual, even when some flows are
nonzero. Box-LP strong duality gives the primal optimum there. Moreover,
equality in the sum of reduced-cost upper bounds forces each nonzero
untied term to use the chosen endpoint, so the remaining tied aggregate
satisfies \eqref{eq:a-wcac-tie-domain}. If every $f_e=0$, any threshold
is feasible with value zero. The argument also covers singleton intervals
and arbitrary ties.
\end{proof}

All endpoints in this lemma are fixed on a sign cell. Both inequalities
in \eqref{eq:a-wcac-tie-domain} are quadratic in $(z,q)$ with a constant
coefficient of $q^2$, and their number does not grow with the number of
tied edges. There are polynomially many realizable sign cells in fixed
$d+1$ dimensions, and at most $|C|$ thresholds per cell.

\begin{lemma}[Complete circulation candidates]
\label{lem:a-wcac-candidates}
The maximum cycle contribution conditional on $z$ is the maximum of
polynomially many feasible candidates, each rational of bounded degree
or of the form $P_2(z)+L_1(z)\sqrt{\Delta(z)}$. Candidate eligibility has
polynomial-size bounded-degree rational semialgebraic descriptions.
These assertions remain true with rational signed arc-flow bounds.
\end{lemma}
\begin{proof}
Fix a closed sign cell and a threshold. Impose $|q|\le R$, its sign
inequalities, \eqref{eq:a-wcac-tie-domain}, and any specified bounds
$\underline x_e\le q+\ell_e(z)\le\overline x_e$. For fixed $z$ the feasible
$q$ set is compact. Each defining polynomial has form
\[
 a q^2+b(z)q+c(z),
\]
with rational constant $a$, affine $b$, and quadratic $c$; linear bounds
are special cases. An optimum of the quadratic $Q_\lambda(z,q)$ occurs
at a boundary, a stationary point, or a feasible point of a flat objective.
A boundary point has at least one active polynomial that is nonconstant
in $q$ after specialization at $z$; otherwise all specialized constraints
would be constant or locally strict there, and it would not be a boundary
point. The artificial bounds $q=\pm R$ are included, so this applies also
to flat objectives.

For a boundary polynomial with $a\ne0$, enumerate both roots
\begin{equation}
 q=\frac{-b(z)\pm\sqrt{b(z)^2-4ac(z)}}{2a}.
 \label{eq:a-wcac-two-roots}
\end{equation}
Substitution in $Q_\lambda$ gives $P_2+L_1\sqrt\Delta$, absorbing the
root sign in $L_1$. The denominator $2a$ is constant. For $a=0$ and
$b(z)\ne0$, substitute $q=-c(z)/b(z)$ and retain the nonzero condition.
On $a=b(z)=0$ the polynomial is constant in $q$ and supplies no boundary
candidate: as the remaining inequalities check, it either excludes the
parameter or imposes no additional $q$ boundary. This handles vanishing
coefficients without dividing by zero.

Write $Q_\lambda=a_Qq^2+b_Q(z)q+c_Q(z)$. If $a_Q\ne0$, include the
stationary candidate $q=-b_Q(z)/(2a_Q)$, whether or not it is a local
maximum; feasibility and later value selection remove harmless extras.
If $a_Q=0$ and $b_Q(z)\ne0$, there is no interior optimum. On
$a_Q=b_Q(z)=0$, include the flat candidate $c_Q(z)$ with the existential
condition that the original $q$ domain is nonempty.

For each nonflat candidate, eliminate $q$ from its defining equation and
all original domain inequalities. The plus root in
\eqref{eq:a-wcac-two-roots} uses the selector $2aq+b\ge0$ and the minus
root uses $2aq+b\le0$; these selectors remain correct when $a<0$ and at
a double root. The rational case uses $b\ne0$. Projection is in $d+1$
variables with bounded degree, so its output size and coefficient lengths
are polynomial. Flat eligibility is the same fixed-dimensional existential
projection without a root equation. Thus no feasible candidate value is
spurious, and the boundary/stationary/flat list includes a local optimum.
\end{proof}

Apply \cref{lem:a-wcac-panels} to every radical candidate, with error
budgets as in \eqref{eq:a-wcac-error}. This time also add to the common
partition the comparison numerators between every pair of surrogate
pieces of the \emph{same} cycle, together with their denominators.
There are polynomially many pieces and pairs, and their degrees and
coefficient lengths remain polynomial. Eligibility and every such
comparison are constant on a sign part. Choose a largest eligible
surrogate for each cycle and then sum, adding bridge endpoint values.
No exact radicals are compared, and no cross-cycle candidate combinations
are enumerated.

For precision accounting, if every candidate of cycle $C$ has error at
most $a_C$, then
\[
 \left|\max_j V_{Cj}-\max_j\widehat V_{Cj}\right|\le a_C.
\]
The candidate selected by the surrogate maximum has true value at least
$\max_j V_{Cj}-2a_C$. Our choice of tolerances ensures
$\sum_C a_C\le a=\eps/16$. Hence the summed surrogate approximates the
conditional joint optimum within $a$, and the actual selected scenario
approximates that surrogate within $a$. Apply the same supremum search
and sampling as before. The interval \eqref{eq:a-wcac-certified-interval}
again encloses the true joint optimum, and the selected physical scenario
at the algebraic sample loses at most $5\eps/16$, including the loss from
selecting candidates by approximate comparisons.

To recover that scenario, recover each selected $q$ from its equation and
root selector over the algebraic sample $z_*$. On a flat branch, sample
one point in its nonempty compact one-dimensional feasible set. For the
tied coordinates, start at the lower end of every interval
\[
 [\min\{L_ef_e,U_ef_e\},\max\{L_ef_e,U_ef_e\}]
\]
and fill the required residual until the sum is $-S$. Divide the resulting
contribution by $f_e$ only if $f_e\ne0$; at zero choose $L_e$. At most one
tied coordinate needs to be interior. This recovers the original
resistance coordinates, not just an eliminated aggregate. In the
directional model give each inactive coefficient any rational value in
its interval. In the symmetric model the recovered active coefficient is
the single shared coefficient; no two independent copies are introduced.
At zero flow choose any allowed original coefficient or pair. These
choices satisfy the original laws.

Process each cycle over the common algebraic $z_*$ and its own local
extension. A nonflat root has degree at most two over the field of $z_*$,
and flat one-dimensional sampling also has a polynomial-size
representation by the fixed-dimensional bounds. Rational operations on
these representations keep polynomial degree and coefficient size. The
output is a list of separate coordinate representations; no primitive
element containing the roots of all cycles is constructed. The resulting
local states share the same nomination and satisfy every cycle equation,
so they form a global physical scenario.

\begin{lemma}[Continuity in the original coefficients]
\label{lem:a-wcac-coefficient-continuity}
Let $\theta$ list the original uncertain coefficients: $m$ entries in the
symmetric model and $2m$ in the directional model. For either model on
any connected graph,
\begin{equation}
 |W(b,\theta)-W(b',\theta')|
 \le K_b\|b-b'\|_1+\|c\|_1B_0^2\|\theta-\theta'\|_1.
 \label{eq:a-wcac-joint-continuity}
\end{equation}
For each terminal potential difference at fixed $b$, its coefficient
change alone is at most $B_0^2\|\theta-\theta'\|_1$.
\end{lemma}
\begin{proof}
For symmetric laws the resistance statement is already
\cref{lem:a-blk-resistance-lipschitz}; the following argument also covers
independent directional coefficients. Smooth by $g_e^\rho(x)=g_e(x)+\rho x$,
with $\rho$ independent of the coefficients. For a terminal objective
$F=\pi_s-\pi_t$, the electrical adjoint solves
$A R^{-1}A^\top h=e_s-e_t$, where
$R=\operatorname{diag}((g_e^\rho)'(x_e))$. Its currents
$j=R^{-1}A^\top h$ have $|j_e|\le1$: nonzero currents follow strictly
decreasing $h$, hence decompose into terminal paths of total weight one.
Differentiating the physical equations at fixed $b$ and pairing with
this adjoint gives
\[
 dF_\rho=\sum_e j_e\,
 \begin{cases}
 x_e|x_e|\,d\beta_e,&\text{symmetric model},\\
 (x_e)_+^2\,d\beta_e^+-(x_e)_-^2\,d\beta_e^-,
                                      &\text{directional model},
 \end{cases}
\]
where $u_+=\max\{u,0\}$ and $u_-=\max\{-u,0\}$. The laws are continuously
differentiable jointly in flow and coefficients, including at zero, and
$R$ is positive, so the differentiation is valid there as well. Every
coefficient derivative has magnitude at most $B_0^2$. Integrate along
the segment inside the original coefficient box and let $\rho\downarrow0$.
The uniform flow and potential bounds and the strict-convexity continuity
argument in \cref{thm:a-wcac-faces} hold uniformly on this compact
coefficient box, so the terminal estimate survives the limit. Normalize
one potential and sum against $|c_v|$. Finally combine this with
\cref{lem:a-wcac-nomination-continuity}, whose $K_b$ is uniform over all
allowed directional coefficients. Sign crossings require no extra case
or nonzero-flow assumption.
\end{proof}

Let $k$ be the length of $\theta$, so $k\le2m$. Choose polynomial-bit
widths $h,\delta>0$ with
$K_bH_1h+\|c\|_1B_0^2k\delta\le\eps/4$.
Round $z_*$ by rational LP inside $P$ and its isolating box, and round each
original coefficient within $\delta$ inside its own rational interval.
In the symmetric model round the one shared coefficient, which preserves
its required equality across the two flow directions. Certified
coordinate isolation takes polynomial time, separately for each local
extension. Singleton intervals are kept exactly. These original
parameters are feasible without maintaining candidate-domain membership,
and the continuity estimate supplies the remaining objective guarantee.
This completes case~2 of \cref{thm:a-wcac-affine}; again negate $c$ for
minimization.

\subsection{Balanced boxes and exact local capacity filters}

\begin{corollary}[Unrestricted nomination boxes at fixed support]
\label{cor:a-wcac-box}
Fix $p$. On a cactus, for a rational balanced nomination box and
$|\supp c|\le p$, both fixed positive asymmetric quadratic laws and joint
independent coefficient intervals in either model of
\cref{thm:a-wcac-affine} admit the value and
rational-scenario guarantees of \cref{thm:a-wcac-affine}. Total cycle rank
is unrestricted. No operating constraints filter this box.
\end{corollary}
\begin{proof}
Apply \cref{thm:a-wcac-faces} with $r=1$. Each of the polynomially many
faces has $O(p)$ free nominations and hence a fixed-dimensional rational
affine parameterization. For joint uncertainty, hold the resistances of a
joint maximizer fixed to identify a face containing an equally good
nomination; the face family is resistance independent. Apply
\cref{thm:a-wcac-affine} on every face with tolerance, for example,
$\eps/8$. The maximum of the certified lower endpoints and the maximum of
the upper endpoints enclose the global optimum. Select the scenario
from a face with the largest lower endpoint. Its optimum is within
$\eps/8$ of the best face optimum and its witness has the stated local
error, which gives the desired global tolerance. Rational disaggregation
reverses all reductions, with arbitrary rational allowed resistances in
removed blocks, and objective preservation gives the original-network
result. The constant-objective case is solved by rational LP. Minimization
follows by changing the sign of $c$.
\end{proof}

\begin{theorem}[Local filters and the rational-output boundary]
\label{thm:a-wcac-capacities}
For the fixed-dimensional affine nomination models of
\cref{thm:a-wcac-affine}, allow rational signed arc bounds
$\underline x_e\le x_e\le\overline x_e$. Feasibility is decidable exactly
in polynomial bit time. If feasible, optimum intervals of width at most
$\eps$ and exactly feasible algebraic $\eps$-optimal original parameters
of polynomial representation size are computable in accuracy-bit
polynomial time.

In the fixed-law model the same algebraic-output conclusion permits a
polynomial-size list of bounded-degree rational polynomial equalities and
inequalities defining a closed set, each involving $z$ and the flows of
one cycle or one bridge. For the joint resistance model the assertion
here is limited to the stated arc bounds.

For either model with arc bounds, let a supplied rational $\sigma>0$
give a nonempty tightened problem
$\underline x_e+\sigma\le x_e\le\overline x_e-\sigma$ on all bounded
ends. A rational scenario feasible for the original bounds and
$\eps$-optimal relative to the tightened optimum can be computed in
accuracy-bit polynomial time, including polynomial dependence on the
encoding length of $\sigma$.
\end{theorem}
\begin{proof}
In the fixed-law model add the local bounds, or local polynomial
constraints, to the chart eligibility formulas before projecting $q$.
The unique physical $q$ and the objective formulas are unchanged.
Substituting the affine flows into a bounded-degree polynomial uses only
$(z,q)$, so the local dimension remains fixed even for long cycles. Bridge
constraints involve $z$ directly. In the joint model the arc bounds are
affine in $(z,q)$ and already covered by \cref{lem:a-wcac-candidates};
arbitrary polynomial filters could add higher-degree circulation
boundaries and are not asserted there.

A global parameter is feasible exactly when every cycle has an eligible
chart or candidate and every bridge bound holds. Take the common
partition with these eligibility conditions and discard parts violating
them; exact real-algebraic feasibility decides whether any part remains.
The parameter and resistance boxes are compact and physical states are
continuous, so the closed filters give a compact feasible scenario
set with attained extrema. The surrogate supremum algorithm and exact
sampling on surviving parts give the stated values and algebraic
parameters. In the joint case use the local tied-coordinate recovery;
all added flow bounds depend only on $(z,q)$, so that recovery preserves
them. No rational LP rounding is used at this stage.

For the tightening claim, the fixed-law nomination flow bound is
\eqref{eq:a-wcac-flow-nomination}. For resistance rounding at fixed $b$,
let $x,y$ be the two flows, with
$\|\theta-\theta'\|_\infty\le\delta$, $k=\dim\theta\le2m$, and
$\beta_L>0$ the smallest allowed directional coefficient. Applying
\cref{lem:a-wcac-coefficient-continuity} to the endpoints of each edge
gives $|\Delta_e-\Delta'_e|\le B_0^2k\delta$. An inverse-law bound gives
\begin{equation}
 \|x-y\|_\infty^2\le
       \frac{2B_0^2(k+1)}{\beta_L}\delta.
 \label{eq:a-wcac-resistance-flow}
\end{equation}
Indeed, for either fixed quadratic law,
$|g_e(u)-g_e(v)|\ge\beta_L|u-v|^2/2$: at equal signs use
$|u^2-v^2|\ge|u-v|^2$, and at opposite signs use
$u^2+v^2\ge(|u|+|v|)^2/2$. Meanwhile changing the original coefficients
within $\delta$ changes $g_e(y_e)$ by at most $B_0^2\delta$, since only
one directional coefficient is active at $y_e$. Subtracting the two
endpoint drops and applying the triangle inequality proves the displayed
bound, even if $x_e,y_e$ have opposite signs; at zero both branches agree.
Combining the nomination and resistance changes bounds the total flow
error by
\[
 \tfrac12 H_1h+
 B_0\sqrt{2(k+1)\delta/\beta_L}.
\]
Choose $h,\delta$ to make this at most $\sigma/2$ and the objective loss
in \eqref{eq:a-wcac-joint-continuity} at most $\eps/4$; in the fixed-law
case omit the resistance term. For example, requiring
$H_1h\le\sigma/2$ and
$\delta\le\beta_L\sigma^2/[32B_0^2(k+1)]$ suffices for the flow condition.
These widths have polynomial bit length. Optimize the tightened
problem with the remaining error budget, then round inside $P$ and the
original resistance intervals. The resulting rational scenario satisfies
the original capacities and loses at most $\eps$ relative to the tightened
optimum. This implies no comparison between the tightened and original
optima, and the rounding need not preserve arbitrary local polynomial
filters from the fixed-law extension.
\end{proof}

\begin{example}[Exact capacities can force irrational nominations]
\label{ex:a-wcac-irrational}
On a coherently oriented four-cycle take $t\in[1,2]$ and
\[
 x=(q+t-1,q-2,q,q-2),\qquad \beta=(1,1,1,2).
\]
Successive flow differences give the balanced rational affine nomination
family $b=(t+1,-t-1,2,-2)$. Impose $|x_3|\le1$, $|x_4|\le1$, and
$|x_1|,|x_2|\le10$. The first two bounds imply both $q\in[-1,1]$ and
$q\in[1,3]$, hence $q=1$. The cycle equation is then
$t^2-1+1-2=t^2-2=0$. Its only feasible parameter is $t=\sqrt2$,
and the other bounds hold. All input data and capacities are rational,
and the symmetric capacities are strictly positive. Thus no rational
parameter is exactly feasible here, so an exactly feasible rational
output cannot be guaranteed, even for fixed resistances.
\end{example}

The box-face corollary does not extend to these filters: pruning,
zero-source contraction, and the box normal-cone argument use the full
unfiltered nomination domain. Likewise, a bound on a potential difference
across many cycles is not a local flow constraint, and exact projection of
such a bound can require comparing a sum of independent radicals.

\subsection{Exact rational resistance profiles at fixed nominations}

\begin{theorem}[Exact rational profile, separately represented value]
\label{thm:a-wcac-exact-profile}
On a cactus with fixed rational balanced nominations, independent
positive rational coefficient intervals in either model of
\cref{thm:a-wcac-affine}, and any rational
balanced objective vector $c$, feasibility under optional rational signed
arc bounds is decidable in polynomial bit time. If feasible, an exactly
maximizing or minimizing rational resistance profile of polynomial
encoding length can be computed in polynomial bit time. Each local
optimal value has algebraic degree at most two; the global value may be
retained as a sum of these separate values. This theorem supplies no
exact scalar comparator for that growing sum.
\end{theorem}
\begin{proof}
Set $d=0$ in \cref{lem:a-wcac-candidates}. All flow-sign, capacity, and
artificial circulation bounds are rational, so their boundary circulations
are rational, as are stationary points of a nonflat rational quadratic
objective. At rational $q$, every $f_e$ is rational, and the tied-aggregate
filling rule recovers rational resistances of polynomial bit length.

An irrational boundary circulation can only come from one of the two
aggregate quadratics in \eqref{eq:a-wcac-tie-domain}. At $S+l_T=0$ the
tied aggregate must attain its minimum, so every nonzero contribution
uses its minimizing resistance endpoint. At $S+u_T=0$ every nonzero tied
contribution analogously uses its maximizing endpoint. Zero-flow
coordinates may use either rational endpoint, and untied coordinates
already use rational endpoints. Thus the recovered original resistance
profile is rational even when its induced $q$ is irrational. In the
directional model choose every inactive coefficient rationally in its
own interval; in the symmetric model retain the shared recovered
coefficient. If both aggregate bounds coincide, either construction works.

For a flat objective choose a boundary point of its nonempty compact
feasible $q$ set, rather than an arbitrary algebraic sample; the
artificial bounds make the set bounded, so such a point exists. At least
one specialized defining polynomial nonconstant in $q$ is active there
(identically zero constraints do not count), which returns to the
rational or aggregate-quadratic boundary cases above. There are
polynomially many candidates on every cycle, each with a rational or
quadratic-algebraic objective value.

Exact comparison of two local values takes polynomial time. One explicit
method writes their difference as $U+V$, with
$U=a+b\sqrt d$ and $V=c\sqrt e$ for rational coefficients and nonnegative
rational radicands. Decide each sign by rational arithmetic and squared
magnitude comparisons. If the signs agree, the answer is immediate. If
they are opposite, the sign is $\sgn(U)\sgn(U^2-V^2)$, and
\[
 U^2-V^2=(a^2+b^2d-c^2e)+2ab\sqrt d
\]
has only one radical. Zero cases are checked first. This recognizes exact
equality without factoring radicands or imposing a separation tolerance;
alternatively, constant-degree root isolation suffices. Select an exact
local winner and retain its rational profile. Bridges check their fixed
flow bounds and choose the objective-maximizing endpoint. The independent
local winners together form a global physical state and an exactly
optimizing rational profile. None of these selections compares the sum of
all block values. Individual radical isolation and rational summation can
enclose that sum to any requested accuracy in polynomial additional bit
time. Negating $c$ yields minimization.
\end{proof}

\begin{example}[A unique interior-resistance optimum]
\label{ex:a-wcac-interior}
On a coherently oriented four-cycle let $b=(0,0,3,-3)$, so
$x=(q,q,q+3,q)$, and let
\[
 \beta\in[1,4]\times[1,3]\times[1,4]\times[1,5],
 \qquad c=(2,-3,-2,3).
\]
The edge weights $w=(2,-1,-3,0)$ satisfy $Aw=c$. Strict monotonicity of
the cycle equation and its values at $q=-3,0$ put every physical
circulation in $(-3,0)$. Subtract $\lambda=-1$ times the cycle equality
from the objective. The three untied reduced costs are strictly negative,
so setting their resistances to their lower bounds gives, for every
feasible profile,
\[
 W\le-4q^2-2(q+3)^2=-6(q+1)^2-12\le-12.
\]
At $q=-1$ the profile $(1,2,1,1)$ gives signed drops $(-1,-2,4,-1)$.
Their sum is zero and their weighted sum is $-12$, so the bound is attained.
Equality forces $q=-1$; then the strictly negative reduced costs force
$\beta_1=\beta_3=\beta_4=1$, and the cycle equation forces $\beta_2=2$.
The optimizer is therefore unique and has an interior second resistance,
so no box corner is optimal. Stationary circulation candidates and tied
resistance recovery cannot be replaced by endpoint-profile enumeration.
\end{example}

The accompanying implementation
\texttt{exact\_weighted\_cactus.py}, in the directory
\texttt{code/potential\_flow\_mpd/},
implements the symmetric-coefficient case of
\cref{thm:a-wcac-exact-profile} for \emph{supplied blocks}:
coherent cycle offsets and weights, resistance intervals, optional local
flow bounds, and fixed-flow bridges. For both extrema it returns exact
rational optimal resistances, separate quadratic encodings of local
circulations and values, and a certified rational interval for the total
value. It does not compute or verify an original graph's block
decomposition, derive its nomination offsets or objective weights,
implement independent directional-coefficient intervals, or implement the
varying-nomination semialgebraic optimizer.
Its exact candidate comparisons handle distinct stored radicands;
arithmetic on local expressions stays within one stored radical, and the
sum across blocks remains separately represented. Its physical-witness,
independent LP-vertex, radical-cancellation, capacity, orientation, and
input-schema checks are regression evidence for this supplied-block
scope; they supplement, but do not replace, the candidate-completeness
proof, and they are not a general untrusted optimality-certificate format.

The structural theorem applies beyond cacti, but the approximation
argument here uses the explicit quadratic, rational, and zero cycle
formulas. It does not give a weighted accuracy-bit algorithm for general
bounded-rank blocks with unbounded total rank. Discrete resistance sets
do not satisfy the tied-interval filling argument, and cross-cycle
resistance correlations do not admit the independent summation step.
